% Статья на конференцию: концепция и алгоритм верхнего уровня.
% Компилировать XeLaTeX. Шаблон конференции не применён — см. README.
\documentclass[12pt,a4paper]{article}

\usepackage{fontspec}
\usepackage{polyglossia}
\usepackage[a4paper,left=2.5cm,right=2.5cm,top=2cm,bottom=2cm]{geometry}
\usepackage{amsmath}
\usepackage{amssymb}
\usepackage{array}
\usepackage{tabularx}
\usepackage{enumitem}
\usepackage{microtype}
\usepackage{caption}
\usepackage[hidelinks]{hyperref}

% Подпись таблицы по-русски: «Таблица 1 — Наименование»
\DeclareCaptionLabelSeparator{dash}{~--- }
\captionsetup{labelsep=dash, justification=raggedright, singlelinecheck=false}

\setdefaultlanguage{russian}
\setotherlanguage{english}
\setmainfont{Times New Roman}
\setlength{\parindent}{1.25cm}
\setlist{nosep, leftmargin=*}
\sloppy

\title{\vspace{-1.5cm}\normalsize УДК 004.4, 004.5, 004.89\\[0.6cm]
\Large\bfseries Семантическое масштабирование при визуальном анализе
дерева поиска методом Монте-Карло}
\author{\normalsize О.\,Д.~Ролдугин, Д.\,В.~Литовкин\\
\normalsize\itshape Волгоградский государственный технический университет}
\date{}

\begin{document}
\maketitle
\thispagestyle{empty}

\noindent\textbf{Аннотация.} Отладка и сравнение вариаций метода поиска по
дереву методом Монте-Карло затруднены размером порождаемого дерева: уже при
секундах на ход число его узлов превышает число элементов, одновременно
различимых на экране. Существующие средства визуализации не сочетают
независимости от конкретной игры с работоспособностью на больших деревьях и
автоматическим управлением детализацией. Предлагается формальная модель, в
которой задача визуального анализа сводится к цепочке из трёх объектов~---
дерева игры, дерева поиска и геометрического дерева~--- и пяти подзадач с
согласованными входами и выходами. Уровень детализации связывается с
масштабом отображения пиксельным критерием видимости, а положение узла
задаётся постоянными нормированными координатами, благодаря чему изменение
детализации не смещает узлы. Связность отображаемого множества и
монотонность его пополнения при приближении следуют из построения.

\smallskip
\noindent\textbf{Ключевые слова:} поиск по дереву методом Монте-Карло,
семантическое масштабирование, визуализация деревьев, уровень детализации,
общее игровое программирование.

\section{Введение}

Поиск по дереву методом Монте-Карло является стандартным подходом к принятию
решений в играх и задачах планирования, а его многочисленные вариации
различаются стратегиями выбора, расширения, симуляции и распространения
оценок~[1]. Понимание поведения конкретной вариации и отладка новой
затруднены: дерево строится несимметрично и стохастически, многие
особенности проявляются на ранних итерациях и маскируются итоговой
статистикой, а размер дерева быстро превосходит возможности экрана. В
наблюдаемых партиях игры <<Четыре в~ряд>> журнал полной партии фиксирует
порядка 42 тысяч итераций, а при десятках секунд на ход счёт итераций
доходит до миллионов.

Положение усугубляется свойствами порождаемых данных. Распределение
посещений резко неоднородно: основную массу поиска сосредоточивает малое
число ветвей, тогда как остальные остаются слаборазвитыми, и скрывать их
нельзя~--- именно они показывают, какие варианты алгоритм отверг. Дерево
переиспользуется между ходами партии, поэтому счётчики посещений предков
накапливаются за всю игру и несопоставимы между ярусами. Наконец, в среде
общего игрового программирования набор игр не фиксирован, поэтому
визуализация оперирует только структурой дерева и статистикой его узлов.
Простое уменьшение масштаба противоречия не разрешает: оно пропорционально
уменьшает все элементы, не различая их значимости. Разрешает его
семантическое масштабирование, при котором изменение масштаба меняет не
только размер элементов, но и состав отображаемой информации~[2,~3].
Настоящая работа формализует применение этого принципа к дереву поиска.

\section{Постановка задачи}

\emph{Дано:} журнал событий поиска~$\sigma$, порождаемый алгоритмом;
момент~$t$, выраженный числом выполненных итераций; параметры камеры~---
видимый домен~$D \subseteq [0,1]$ и ширина окна~$W$ в пикселях; бюджет
числа одновременно изображаемых узлов~$B$.

\emph{Требуется:} построить изображение состояния поиска на момент~$t$, в
котором уровень детализации согласован с масштабом, положение узла не
зависит от того, какая часть дерева загружена, а скрытые поддеревья
представлены обобщёнными участками, сохраняющими их статистику.

\section{Формальная модель}

\subsection{Дерево игры}

Игра задаётся набором $G = \langle R, S, s_0, A, terminal, legal, next,
payoff \rangle$. Вершиной порождаемого ею дерева~$T_{game}$ служит
\emph{история}~--- последовательность совместных действий $h = \langle a_1,
\dots, a_k \rangle$, выполнимая из~$s_0$; дуги ведут в продолжения $h \cdot
a$ при $a \in legal(state(h))$. Вершиной объявлена история, а не состояние,
намеренно: одно состояние достижимо разными последовательностями ходов, и
отождествление таких вершин превратило бы дерево в ориентированный
ациклический граф. Дерево игры \emph{дано} правилами и не строится.

\subsection{Дерево поиска}

Поиск порождает журнал из событий трёх видов: границы хода, расширения
вершины с перечнем созданных потомков и итерации с указанием выбранного
листа, вершины симуляции, полученных платежей и глубины доигрывания. Пути
этапов выбора и обратного распространения не хранятся: оба проходят по
единственному пути от корня к вершине симуляции и выводятся из цепочки
предков. Дерево поиска на момент~$t$ есть

\begin{equation}
T_{search}(t) = (V_t,\; E_t,\; N_t,\; Q_t), \qquad V_t \subseteq V_{game},
\end{equation}

где $N_t$~--- число посещений, $Q_t$~--- средняя накопленная оценка. Модель
опирается только на перечисленные события, поэтому любая вариация метода,
порождающая такой журнал, укладывается в неё без изменений.

\subsection{Геометрическое дерево}

Каждой вершине назначается интервал нормированной горизонтальной оси и ярус
её глубины: $T_{geom}(t) = (V_t, x_t, y)$, где $x_t(v) = [x_0(v), x_1(v)]
\subseteq [0,1]$ и $y(v) = depth(v)$. Геометрия подчинена двум требованиям:

\begin{equation}
x_t(c) \subseteq x_t(p) \quad \forall (p,c) \in E_t,
\qquad \sum_{c \in children(p)} w(c) = w(p),
\label{eq:geom}
\end{equation}

где $w(v) = x_1(v) - x_0(v)$. Вложенность кодирует отношение <<предок~---
потомок>> положением прямоугольников, поэтому дуги не изображаются вовсе, а
разбиение без остатка означает, что потомки делят полосу предка полностью.

Доли потомков вычисляются среди соседей по общему предку, а не относительно
корня: при переиспользовании дерева абсолютные счётчики несопоставимы между
ярусами. Ширина и значимость разводятся по двум каналам: ширина получает
нижний порог $\delta = \min(0{,}02;\, 0{,}5/n)$, чтобы раскрытые, но ещё не
посещённые ходы оставались видимы, а значимость порога не имеет и служит
основой цветового кодирования.

При ширине окна~$W$ и домене~$D$ плотность отображения равна $\rho = W /
(d_1 - d_0)$, и вершина считается допустимой к изображению, если

\begin{equation}
x(v) \cap D \ne \varnothing \;\wedge\; w(v)\cdot\rho \ge s_{min},
\label{eq:adm}
\end{equation}

где $s_{min}$~--- перцептивный минимум экранной ширины. Условие~(\ref{eq:adm})
и делает масштабирование семантическим: при уменьшении домена плотность
растёт, и порог проходят всё более мелкие узлы.

\section{Алгоритм верхнего уровня}

Задача решается последовательным применением пяти подзадач; выход каждой
служит входом следующей (таблица~\ref{tab:io}).

\smallskip
\noindent\textbf{Алгоритм 1.} \emph{Вход:} $\sigma$, $t$, $D$, $W$, $B$.
\emph{Выход:} изображаемые узлы~$R$ и агрегаты скрытого~$H$.
\begin{enumerate}[label=\arabic*.]
  \item $T \leftarrow \textsc{Реконструировать}(\sigma, t)$;
  \item $T^{*} \leftarrow \textsc{ОпорноеСостояние}(\sigma, t)$;
  \item $g \leftarrow \textsc{ПостроитьГеометрию}(T, T^{*})$;
  \item $R \leftarrow \textsc{ОтобратьВидимые}(g, D, W/(d_1-d_0), B)$;
  \item $H \leftarrow \textsc{СвернутьСкрытое}(g, R)$.
\end{enumerate}

\begin{table}[t]
\centering
\small
\caption{Входы и выходы подзадач, согласованные по цепочке}
\label{tab:io}
\begin{tabularx}{\linewidth}{>{\raggedright\arraybackslash}p{0.30\linewidth}>{\raggedright\arraybackslash}X>{\raggedright\arraybackslash}p{0.26\linewidth}}
\hline
Подзадача & Вход & Выход \\
\hline
Реконструкция & журнал, момент & дерево поиска \\
Опорное состояние & журнал, границы хода & опорное дерево \\
Построение геометрии & дерево поиска, опорное дерево & геометрическое дерево \\
Отбор видимых & геометрия, камера, бюджет & изображаемое множество \\
Свёртка скрытого & геометрия, изображаемое множество & агрегаты со сводками \\
\hline
\end{tabularx}
\end{table}


\smallskip
Реконструкция применяет события журнала по порядку: расширение подвешивает
потомков, итерация увеличивает счётчики на пути от вершины симуляции к
корню, граница хода меняет корень. Геометрия назначается одним обходом в
глубину по опорному состоянию. Отбор ведётся обходом от корня по убыванию
ширины с проверкой~(\ref{eq:adm}) и остановкой по бюджету. Свёртка строит
агрегат на каждую серию подряд идущих неизображаемых потомков и снабжает
его сводкой: число свёрнутых ходов и узлов, суммарные посещения и их доля,
глубина, диапазон оценок.

Опорное состояние~$T^{*}$ разделяет два режима просмотра. Для записанного
хода им служит итоговое состояние: внутри хода узлы только добавляются,
поэтому $V_t \subseteq V_{T_h}$, геометрия неподвижна весь ход и узел
появляется сразу на своём окончательном месте. При наблюдении в реальном
времени итогового состояния ещё не существует, опорным служит текущее
дерево, и устойчивость не гарантируется.

\section{Свойства решения}

Из построения следуют три свойства, которые не требуют отдельных проверок.

\emph{Связность.} Из вложенности~(\ref{eq:geom}) следует $w(p) \ge w(c)$,
поэтому допустимость потомка по~(\ref{eq:adm}) влечёт допустимость предка, а
в обходе потомок рассматривается лишь после включения предка. Отображаемое
множество всегда связно, и отдельная процедура достройки предков не нужна.

\emph{Монотонность при приближении.} Порог в~(\ref{eq:adm}) монотонен по
ширине, поэтому при росте~$\rho$ вновь допустимые узлы заведомо мельче ранее
включённых и не вытесняют их из бюджета: приближение только добавляет
подробности.

\emph{Сохранение массы.} Из разбиения~(\ref{eq:geom}) следует, что на каждом
ярусе сумма ширин изображаемых узлов и агрегатов равна ширине полосы
предка. Пользователь всегда видит всю статистику поиска, меняется лишь её
гранулярность.

Измерения на реализации подтверждают расчётные свойства: на деревьях до
193~тысяч узлов размер отображаемого множества не превышает 460 узлов при
бюджете 600 и задержке ответа не более 3~мс, а смещение узлов при изменении
детализации равно нулю против сотен тысяч пикселей у раскладок с пересчётом
при каждой догрузке.

\section{Заключение}

Предложена формальная модель визуального анализа дерева поиска, в которой
задача выводится из трёх объектов и раскладывается на пять подзадач с
согласованными входами и выходами. Уровень детализации связан с масштабом
пиксельным критерием видимости, положение узла задано постоянными
нормированными координатами, а скрытые поддеревья сохранены как агрегаты со
статистикой. Связность, монотонность и сохранение массы следуют из
построения, а не обеспечиваются дополнительными механизмами. Дальнейшая
работа~--- учёт нераскрытых продолжений при вариациях с прогрессивным
расширением и оценка вклада отдельных решений.

\begin{thebibliography}{9}\footnotesize
\setlength{\itemsep}{0pt}\setlength{\parskip}{0pt}
\bibitem{} Świechowski~M., Godlewski~K., Sawicki~B., Mańdziuk~J. Monte Carlo
Tree Search: A Review of Recent Modifications and Applications~//
Artificial Intelligence Review. 2023. Vol.~56. P.~2497--2562.
\bibitem{} Perlin~K., Fox~D. Pad: An Alternative Approach to the Computer
Interface~// Proceedings of SIGGRAPH~'93. 1993. P.~57--64.
\bibitem{} Cockburn~A., Karlson~A., Bederson~B.~B. A Review of
Overview+Detail, Zooming, and Focus+Context Interfaces~// ACM Computing
Surveys. 2009. Vol.~41, №~1. P.~1--31.
\bibitem{} Gregg~B. The Flame Graph~// Communications of the ACM. 2016.
Vol.~59, №~6. P.~48--57.
\bibitem{} Coulom~R. Treemaps for Search-Tree Visualization~// Proceedings
of the 7th Computer Olympiad Computer-Games Workshop. Maastricht, 2002.
\bibitem{} Elmqvist~N., Fekete~J.-D. Hierarchical Aggregation for
Information Visualization: Overview, Techniques, and Design Guidelines~//
IEEE Transactions on Visualization and Computer Graphics. 2010. Vol.~16,
№~3. P.~439--454.
\bibitem{} Genesereth~M., Love~N., Pell~B. General Game Playing: Overview of
the AAAI Competition~// AI Magazine. 2005. Vol.~26, №~2. P.~62--72.
\bibitem{} Литовкин~Д.~В., Ролдугин~О.~Д., Якимов~Г.~А., Николаенко~С.~Д.,
Смутин~Д.~А. Динамическая визуализация построения MCTS-дерева в среде GGP
Base Package~// Моделирование, оптимизация и информационные технологии.
2026. Т.~14, №~7. С.~1--11.
\end{thebibliography}

\end{document}
